\documentclass[../../../include/open-logic-section]{subfiles}

\begin{document}
\olfileid[hi]{sth}{ordinals}{zfm}	
\olsection{$\ZFminus$: एक महत्त्वपूर्ण पड़ाव}

प्रतिस्थापन का औचित्य कैसे दिया जाए (यदि दिया भी जा सके), यह प्रश्न सरल
नहीं। इसलिए इसे हम \olref[replacement][]{chap} के लिए सुरक्षित रखेंगे।
किंतु प्रतिस्थापन जोड़ने पर हम एक और महत्त्वपूर्ण पड़ाव पर पहुँच गए हैं। अब
हमारे पास सिद्धांत $\ZFminus$ के लिए आवश्यक सभी अभिगृहीत हैं। विस्तार से:

\begin{defn}
सिद्धांत $\ZFminus$ के अभिगृहीत ये हैं: विस्तारिता, सम्मिलन, युग्म, घात
समुच्चय, अनंतता, और पृथक्करण तथा प्रतिस्थापन योजनाओं के सभी दृष्टांत। दूसरे
शब्दों में, $\ZFminus$, $\Zminus$ में प्रतिस्थापन जोड़ता है।
\end{defn}
\noindent
यह \emph{ज़ेर्मेलो--फ़्रेंकेल} समुच्चय सिद्धांत को दर्शाता है (उसमें से कुछ
\emph{घटाकर}, जिस पर हम बाद में आएँगे)। फ़्रेंकेल को यह सम्मान मिलता है,
क्योंकि \citeyear{Fraenkel1922} में प्रतिस्थापन के सूत्रीकरण का श्रेय उन्हें
दिया जाता है, यद्यपि पहला सटीक सूत्रीकरण \citet{Skolem1922} ने दिया था।

\end{document}
